\documentclass[10pt,a4paper]{amsart}
\usepackage[margin=19mm]{geometry}
\usepackage{amsmath,amssymb,mathtools}
\usepackage[scr=rsfs,cal=euler]{mathalfa}
\usepackage{xfrac}
\usepackage[unicode,hidelinks,bookmarks=false]{hyperref}
\newcommand{\ie}{i.e.\ }
\newcommand{\cf}{cf.\ }
\newcommand{\resp}{respectively\ }
\newcommand{\Subsection}{\subsection}
\providecommand{\nobreakdash}{\nobreak\hbox{-}}
\setlength{\emergencystretch}{2em}
\multlinegap0pt
\usepackage{bm}
\usepackage{bbm}
\usepackage{dsfont}
\usepackage{xspace}
\usepackage{cancel}
\usepackage{mathtools}
\usepackage{amsthm}
\makeatletter
\@ifundefined{thm}{\newtheorem{thm}{Thm}}{}
\@ifundefined{lem}{\newtheorem{lem}[thm]{Lemma}}{}
\@ifundefined{prop}{\newtheorem{prop}[thm]{Proposition}}{}
\@ifundefined{rem}{\newtheorem{rem}[thm]{Remark}}{}
\makeatother
\newcommand{\E}{\mathcal{E}}
\newcommand{\R}{\mathbb{R}}
\newcommand{\sym}{\text{sym}}
\renewcommand{\S}{{\mathcal{S}_n}}
\title[Very weak solutions of the Dirichlet problem for 2-Hessian e]{Very weak solutions of the Dirichlet problem for 2-Hessian equation}
\author{Tongtong Li, Guohuan Qiu}
\date{Source: arXiv:2403.09356v1 (CC BY 4.0)}
\begin{document}
\maketitle

\noindent\emph{Source and licence.} Very weak solutions of the Dirichlet problem for 2-Hessian equation. Version arXiv:2403.09356v1 (CC BY 4.0), distributed under the Creative Commons Attribution 4.0 International licence, as recorded in the arXiv licence metadata for this exact version and shown on the abstract page https://arxiv.org/abs/2403.09356v1. Full rights notice: licenses/arxiv-2403.09356-CC-BY-4.0.txt.\par\smallskip

\noindent\textit{Editorial scope.} Counted unit: the Prop whose source label is \texttt{prop1} in the article (source0.tex lines 585--598, source label \texttt{prop1}), delivered together with the article's own complete original proof of it (source0.tex lines 599--783, one \texttt{proof} environment of 185 lines). No mathematics is written, repaired or re-derived: statement and proof are the article's own text line for line. Required context retained uncounted: bounds (lines 234--239); phil1 (lines 251--255); phil2 (lines 251--255); phil3 (lines 251--255); 1 (lines 261--263); 4 (lines 293--296); lem2 (lines 469--478); 881 (lines 581--583). Every definition, display and label of those lines is retained unchanged. Omitted from this package and not redistributed: 1--233, 240--250, 256--260, 264--292, 297--468, 479--580, 584--584, 784--1463. Nothing inside a retained excerpt is omitted or altered: every retained body is delivered exactly as the article's lines are written, so no omission receipt of any kind accompanies this package. Citations: the retained excerpts carry no citation command at all, so no bibliography entry of the article is transcribed and no reference is redistributed. Reference adaptation: none was needed and none was made; every reference inside the delivered unit resolves inside it, so no unresolved reference or citation is produced. Printed slips kept literal: no printed word, symbol, hypothesis or number of the counted statement or of its proof was altered. Layout and class: the article's own class is \texttt{amsart} with packages including inputenc, graphicx, amsmath, amssymb, amsthm, mathrsfs, color, times, textcomp, verbatim, yfonts, geometry, hyperref, mathtools; the wrapper is the lane's pinned \texttt{amsart} editorial wrapper at 10pt on A4, which restates the theorem-like environments the retained text needs and adds the article's own macro definitions that the retained text needs (\texttt{\textbackslash E}, \texttt{\textbackslash R}, \texttt{\textbackslash S}, \texttt{\textbackslash sym}), with extra wrapper packages (bm, bbm, dsfont, xspace, cancel, mathtools). The delivered numbering is the unit's own. Assets: no retained span contains a figure environment, a graphics-inclusion command, a PDF-page inclusion, a table or a TikZ picture; no image file is copied into this package, the package declares no asset dependency and no image file is redistributed with it. Licence: version arXiv:2403.09356v1 of this article is distributed under the Creative Commons Attribution 4.0 International licence, as recorded in the arXiv licence metadata for this exact version and shown on the abstract page https://arxiv.org/abs/2403.09356v1. A full rights notice is delivered at licenses/arxiv-2403.09356-CC-BY-4.0.txt. Characters: every retained excerpt is pure ASCII, so the delivered engine, XeLaTeX with its default Latin Modern text font, reports no missing character.
\begin{equation}\label{bounds}
\begin{aligned}
    |\partial_t^k\Gamma_1(s,t)|\le Cs,\quad |\partial_s\partial_t^k\Gamma_1(s,t)|\le C,\\
    |\partial_t^k\Gamma_2(s,t)|\le Cs^2,\quad |\partial_s\partial_t^k\Gamma_2(s,t)|\le Cs,\quad |\partial_s^2\partial_t^k\Gamma_2(s,t)|\le C.
\end{aligned}
\end{equation}

\begin{gather}
    [h*\phi_l]_{r}\le [h]_{r}, \label{phil3}\\
    [h-h*\phi_l]_{r}\le C[h]_{r+s}l^{s},\quad [h*\phi_l]_{r+s}\le C[h]_r l^{-s},\label{phil1}\\
    [(h_1h_2)*\phi_l-(h_1*\phi_l)(h_2*\phi_l)]_r\le C\|h_1\|_1\|h_2\|_1 l^{2-r},\label{phil2}
\end{gather}

\begin{equation}\label{1}
    \sigma_2(\nabla^2 v) = f \quad\text{in }\Omega.
\end{equation}

\begin{align}
    -L(A)=f\quad\text{in}\ \Omega, \label{2}\\
    A=\frac{1}{2}\nabla v\otimes\nabla v+\sym\nabla w\quad\text{in}\ \Omega.\label{4}
\end{align}

\begin{lem}\label{lem2}
    Let $N_*=\frac{1}{2}n(n+1)$ be the dimension of $\S$. Then there exist constants $\sigma_*\in(0,1/2), C_*>c_*>0$, unit vectors $ \xi_1,..., \xi_{N_*}$ in $\R^n$ and $C^\infty$ functions $d_1,...,d_{N_*}$ from $\S$ to $\R$, depending only on $n$, such that for any $D\in\S$ and
    \begin{equation*}
	\|D-Id\|_0\le\sigma_*,
    \end{equation*}
    we have
    \begin{equation*}
        D=\sum_{i=1}^{N_*}d_i^2 \xi_i\otimes \xi_i,\ d_i=d_i(D)\in[c_*,C_*],\ \forall i.
    \end{equation*}
\end{lem}

\begin{equation}\label{881}
   D_q:= A-\frac{1}{2}\nabla V_q\otimes\nabla V_q-\sym\nabla W_q-\delta_{q+1}Id.
\end{equation}

\begin{prop}\label{prop1}
    There exist some positive universal constants $\sigma\in(0,\sigma_*/3)$ and $K>1$, such that: if $(V_q,W_q)\in C^2(\bar\Omega)\times C^2(\bar\Omega,\R^n)$ satisfy
    \begin{gather}
         \|(V_q,W_q)\|_1\le \sqrt{K}, \label{prop1.a1}\\
         \|(V_q,W_q)\|_2\le K\delta_q^\frac{1}{2}\lambda_q, \label{prop1.a2}\\
         \|D_q\|_0\le{\sigma}\delta_{q+1}. \label{prop1.a3}
    \end{gather}
    we can construct $(V_{q+1},W_{q+1})\in C^2(\bar\Omega)\times C^2(\bar\Omega,\R^n)$ with
    \begin{gather}
        \|(V_{q+1}-V_q, W_{q+1}-W_q)\|_1\le K\delta_{q+1}^\frac{1}{2}, \label{prop1.0}\\
        \|(V_{q+1},W_{q+1})\|_2\le K\delta_{q+1}^\frac{1}{2}\lambda_{q+1},\label{prop1.1}\\
         \|D_{q+1}\|_0\le{\sigma}\delta_{q+2}.\label{prop1.2}
    \end{gather}
\end{prop} 

\begin{proof}
In the proof all implicit (universal) constants are independent of $\sigma,K$. For convenience we write
\begin{equation}\label{mu0}
    \mu_0:=K\delta_{q+1}^{-\frac{1}{2}}\delta_q^\frac{1}{2}\lambda_q,
\end{equation}
thus
\begin{equation}\label{prop1.a2'}
    \|(V_q,W_q)\|_2\le \delta_{q+1}^\frac{1}{2}\mu_0.
\end{equation}
Let $\tilde A=A*\phi_l,\tilde v=V_q*\phi_l,\tilde w=W_q*\phi_l$ be the mollification of $A,V_q,W_q$, respectively, at length $l$. 
\begin{rem}
    Unlike before, here we need mollification to avoid the dependence on uncontrolled higher derivatives of $(V_q,W_q)$, see Remark \ref{rem998}. Mollification is usually done in the entire space $\R^n$. To this end we actually first extend $v^b,f$ to $C_c(\tilde\Omega)$ for some $\tilde\Omega\Supset\Omega$. and then solve \eqref{1} in $\tilde\Omega$. This does not affect the arguments in the next section, since we cut-off everything there inside $\Omega$.

\end{rem}
Denote
\begin{equation*}
    \tilde D := \tilde A-\frac{1}{2} \nabla\tilde v\otimes\nabla\tilde v-\sym\nabla\tilde w-\delta_{q+2}Id.
\end{equation*}
Then by \eqref{881}
\begin{equation*}
    \tilde D-\delta_{q+1}Id= D_q*\phi_l+\frac{1}{2}[(\nabla V_q\otimes\nabla V_q)*\phi_l-\nabla\tilde v\otimes\nabla\tilde v]-\delta_{q+2}Id.
\end{equation*}
Assume
\begin{equation}\label{ineq02}
    \delta_{q+2}\le {\sigma}\delta_{q+1},
\end{equation}
and injecting \eqref{phil3} \eqref{phil2} and assumptions \eqref{prop1.a1} \eqref{prop1.a2'} \eqref{prop1.a3}, we see
\begin{align*}
    \|\tilde D-\delta_{q+1}Id\|_0
    &\le \|D_q*\phi_l\|_0+\frac{1}{2}\|(\nabla V_q\otimes\nabla V_q)*\phi_l-\nabla\tilde v\otimes\nabla\tilde v\|_0+\delta_{q+2}\\
    &\le \|D_q\|_0 + Cl^2\|V_q\|_2^2 + \delta_{q+2}\\
    &\le 2{\sigma}\delta_{q+1}+C(\mu_0 l)^2\delta_{q+1}
\end{align*}
Now we take
\begin{equation}\label{l}
    l:=\frac{\sigma}{C\mu_0} = \frac{\sigma}{CK}\delta_{q+1}^\frac{1}{2}\delta_q^{-\frac{1}{2}}\lambda_q^{-1}
\end{equation}
for some constant $C>1$ large enough, independent of $\sigma$ and $K$,  
% \begin{rem}\color{red}
%     We can also take $l$ a bit smaller and put this term into $D_{q+1}$, similar for $A-\tilde A$. 
% \end{rem}
such that
\begin{equation*}
    \|\tilde D-\delta_{q+1}Id\|_0\le 3\sigma\delta_{q+1},
    \quad\text{i.e.}\quad
    \|\frac{\tilde D}{\delta_{q+1}}-Id\|_0\le 3\sigma\le\sigma_*.
\end{equation*}
Similarly 
\begin{equation*}
    \|\tilde D-\delta_{q+1}Id\|_k \le C\sigma\delta_{q+1}l^{-k},\quad\forall 1\le k\le 3.
\end{equation*}
Now Lemma \ref{lem2} applies to $\frac{\tilde D}{\delta_{q+1}}$, yielding
\begin{equation}\label{882}
    \tilde D=\sum_{i=1}^{N_*}a_i^2\xi_i\otimes\xi_i,\quad\text{i.e.}\quad
    \tilde A-\frac{1}{2} \nabla\tilde v\otimes\nabla\tilde v-\sym\nabla\tilde w-\delta_{q+2}Id=\sum_{i=1}^{N_*}a_i^2\xi_i\otimes\xi_i,
\end{equation}
where
\begin{gather}
    \|a_i\|_0\lesssim\delta_{q+1}^\frac{1}{2}, \label{721}\\
    \|\nabla^k a_i\|_0\lesssim\delta_{q+1}^\frac{1}{2}\|\frac{\tilde D}{\delta_{q+1}}-Id\|_k\lesssim \sigma\delta_{q+1}^\frac{1}{2}l^{-k},\ \forall 1\le k\le 3. \label{722}
\end{gather}
Now we define, for $1\le i\le N_*$
\begin{align}
    &(v_0, w_0)=(\tilde v, \tilde w)\\
    & v_i= v_{i-1}+\frac{1}{\mu_i}\Gamma_1(a_i, \mu_i x\cdot \xi_i),\\
    & w_i= w_{i-1}-\frac{1}{\mu_i}\Gamma_1(a_i,\mu_i x\cdot \xi_i)\nabla v_{i-1}+\frac{1}{\mu_i}\Gamma_2(a_i,\mu_i x\cdot \xi_i) \xi_i,\\
    & (V_{q+1},W_{q+1})=(v_{N_*},w_{N_*}),
\end{align}
with constants \begin{equation*}
\mu_0\le\mu_1\le...\le\mu_{N_*}=\lambda_{q+1}
\end{equation*}
determined later. Using \eqref{phil1} \eqref{phil2} and \eqref{prop1.a2'}, we see
\begin{gather}
    \|(v_0-V_q,w_0-W_q)\|_1\lesssim \|(V_q.W_q)\|_2 l\lesssim \delta_{q+1}^\frac{1}{2}\mu_0 l\lesssim \sigma\delta_{q+1}^\frac{1}{2},\label{401}\\
    \|v_0\|_{2+k}\lesssim \|V_q\|_2 l^{-k}\lesssim \delta_{q+1}^\frac{1}{2}\mu_0 l^{-k},\ k=0,1. \label{402}
\end{gather}
We first estimate $v_i$ as follows: by  \eqref{bounds} and \eqref{721} \eqref{722}, for $0\le k\le 3$ there holds
\begin{equation}\label{403}
\begin{aligned}
    \|v_i- v_{i-1}\|_k
    &\lesssim \frac{1}{\mu_i}(\mu_i^k\|a_i\|_0+\mu_i^{k-1}\|a_i\|_1+...+\|a_i\|_k)\\
    &\lesssim\frac{1}{\mu_i}(\mu_i^k\delta_{q+1}^\frac{1}{2}+\mu_i^{k-1}\cdot\sigma\delta_{q+1}^\frac{1}{2}l^{-1}+...+\sigma\delta_{q+1}^\frac{1}{2}l^{-k})\\
    &\lesssim \delta_{q+1}^\frac{1}{2}\mu_i^{k-1}
\end{aligned}
\end{equation}
where we assumed $l^{-1}\le\mu_1$, namely
\begin{equation}\label{ineq03}
    \frac{C\mu_0}{\sigma}\le\mu_1\le...\le\mu_{N_*}=\lambda_{q+1}.
\end{equation}
As a byproduct
\begin{gather*}
    \|\nabla v_i\|_0\le \|\nabla V_q\|_0+C\delta_{q+1}^\frac{1}{2}\le \sqrt{K}+C\lesssim\sqrt{K},\\
    \|\nabla^2 v_i\|_0\lesssim \delta_{q+1}^\frac{1}{2}(\mu_0+\mu_1+...+\mu_i)\lesssim\delta_{q+1}^\frac{1}{2}\mu_i,\\
    \|\nabla^3 v_0\|_0\lesssim \delta_{q+1}^\frac{1}{2}\mu_0 l^{-1},\\
    \|\nabla^3 v_i\|_0\lesssim \delta_{q+1}^\frac{1}{2}\mu_0 l^{-1}+\delta_{q+1}^\frac{1}{2}(\mu_1^2+...+\mu_i^2)\lesssim \delta_{q+1}^\frac{1}{2}\mu_i^2.
\end{gather*}
These will be used to estimate $w_i$ and $\E_i$.

Now turn to $w_i$. By \eqref{bounds} \eqref{721} \eqref{722} and above bounds for $v_{i-1}$,  we have
\begin{equation}\label{404}
\begin{aligned}
    \|w_i-w_{i-1}\|_1
    &\lesssim \frac{1}{\mu_i}(\|a_i\|_0\|\nabla v_{i-1}\|_0\mu_i+\|a_i\|_1\|\nabla v_{i-1}\|_0+\|a_i\|_0\|\nabla v_{i-1}\|_1)\\
    &\quad +\frac{1}{\mu_i}(\|a_i^2\|_0\mu_i+\|a_i^2\|_1)\\
    &\lesssim  \frac{1}{\mu_i}(\delta_{q+1}^\frac{1}{2} \cdot\sqrt{K}\mu_i+\delta_{q+1}^\frac{1}{2}l^{-1}\cdot\sqrt{K}+\delta_{q+1}^\frac{1}{2}\cdot \delta_{q+1}^\frac{1}{2}\mu_{i-1})\\
    &\quad +\frac{1}{\mu_i}(\delta_{q+1}\mu_i+\delta_{q+1}l^{-1})\\
    &\lesssim\sqrt{K}\delta_{q+1}^\frac{1}{2}.
\end{aligned}
\end{equation}
Similarly 
\begin{equation}\label{405}
    \|w_i-w_{i-1}\|_2\lesssim \sqrt{K}\delta_{q+1}^\frac{1}{2}\mu_i.
\end{equation}
\begin{rem}\label{rem998}
Although to solve \eqref{4} only requires $W_q$ converges in $C^1$ sense, \eqref{405} is necessary: note that to estimate $\|\nabla^2 a_i\|_0$ used in \eqref{403}, we need to mollify $W_q$, which produces an error $\|w_0-W_q\|_1$, hence the bound of $\|W_q\|_2$, in particular \eqref{405}, is unavoidable. 
\end{rem}
By \eqref{401} \eqref{403} \eqref{404}
\begin{equation}\label{K6}
\begin{aligned}
    &\|(V_{q+1}-V_q, W_{q+1}-W_q)\|_1\\
    &\lesssim \|(v_0,w_0)-(V_q,W_q)\|_1+\sum_{i=1}^{N_*}(\|v_i-v_{i-1}\|_1+\|w_i-w_{i-1}\|_1)\\
    &\lesssim\sqrt{K}\delta_{q+1}^\frac{1}{2}
    \le K\delta_{q+1}^\frac{1}{2},
\end{aligned}
\end{equation}
and similarly by \eqref{402} \eqref{403} \eqref{405}
\begin{equation}\label{K2}
    \|(V_{q+1}-V_q, W_{q+1}-W_q)\|_2\lesssim \sqrt{K}\delta_{q+1}^\frac{1}{2}(\mu_0+...+\mu_{N_*})
    \le K\delta_{q+1}^\frac{1}{2}\lambda_{q+1},
\end{equation}
as long as $K>1$ is large enough. This verifies \eqref{prop1.0} \eqref{prop1.1}. 

It remains to prove \eqref{prop1.2}. The $i$-th step error is
\begin{align*}
    {\E}_i
    &={a}_i^2 \xi_i\otimes \xi_i
    +(\frac{1}{2}\nabla  {v}_{i-1}\otimes\nabla  {v}_{i-1}+\sym\nabla   {w}_{i-1})
    -(\frac{1}{2}\nabla  {v}_i\otimes\nabla  {v}_i+\sym\nabla   {w}_i)\\
    & =\frac{1}{ {\mu}_i}\Gamma_1\nabla^2  {v}_{i-1}
    -\frac{1}{ {\mu}_i}(\partial_{s}\Gamma_2+(\partial_{s}\Gamma_1)(\partial_{t}\Gamma_1))\sym(\nabla {a}_i\otimes \xi_i)\\
    & -\frac{1}{2 {\mu}_i^2}|\partial_{s}\Gamma_1|^2\nabla {a}_i\otimes\nabla {a}_i.
\end{align*}
Then by \eqref{881} \eqref{882}
\begin{equation}
    D_{q+1}=A-\tilde A+\sum_{i=1}^{N_*}\E_i.
\end{equation}
Recalling \eqref{bounds}, one obtains
\begin{equation}
\begin{aligned}
    \|\E_i\|_0
    &\lesssim\frac{1}{\mu_i}\|\Gamma_1\|_0\|\nabla^2 v_i\|_0 +\frac{1}{\mu_i}(\|\partial_s\Gamma_2\|_0+\|\partial_s\Gamma_1\|_0\|\partial_t\Gamma_1\|_0)\|\nabla a_i\|_0\\
    &+\frac{1}{\mu_i^2}\|\partial_s\Gamma_1\|_0^2\|\nabla a_i\|_0^2\\
    &\lesssim\frac{1}{\mu_i}\delta_{q+1}^\frac{1}{2}\cdot \delta_{q+1}^\frac{1}{2}\mu_{i-1}
    + \frac{1}{\mu_i}\cdot\delta_{q+1}^\frac{1}{2}\cdot\sigma\delta_{q+1}^\frac{1}{2}l^{-1}
    +\frac{1}{\mu_i^2}(\sigma\delta_{q+1}^\frac{1}{2}l^{-1})^2\\
    &\lesssim \delta_{q+1}\frac{\mu_{i-1}}{\mu_i}
\end{aligned}
\end{equation}
Besides, since $A\in C^{2,\kappa}(\bar\Omega,\S)$,
\begin{equation*}
    \|A-\tilde A\|_0\lesssim \|A\|_1 l\lesssim l.
\end{equation*}
So we conclude
\begin{equation*}
    \|D_{q+1}\|_0\lesssim l+\delta_{q+1}(\frac{\mu_0}{\mu_1}
    +\frac{\mu_1}{\mu_2}+...+\frac{\mu_{N_*-1}}{\mu_{N_*}}).
\end{equation*}
We now take frequencies $\mu_i$ as
\begin{equation*}
    \mu_i:=\mu_0^{1-\frac{i}{N_*}}\mu_{N_*}^\frac{i}{N_*}
\end{equation*}
and require
\begin{equation}\label{ineq**}
    l\le \frac{\sigma\delta_{q+2}}{C},\quad 
    \frac{\mu_0}{\mu_{N_*}}\le
    (\frac{\sigma\delta_{q+2}}{C\delta_{q+1}})^{N_*}
\end{equation}
for some universal constant $C>1$ large enough. 
In particular \eqref{ineq03} holds.
Now we get
\begin{equation*}
\|D_{q+1}\|_0\lesssim l+\delta_{q+1}(\frac{\mu_0}{\mu_{N_*}})^\frac{1}{N_*}\le{\sigma}\delta_{q+2}
\end{equation*}
which is exactly \eqref{prop1.2}. This concludes the proof of Proposition \ref{prop1}.
\end{proof}

\end{document}
